\section{Labels and \bsq{break}/\bsq{continue} Statements}\label{sec:LabelsAndBreakStatements}
According to \autocite{ORACLE_DOC_LANGUAGE_SPECIFICATION:LabeledStatements}, in Java a label\index{label} can be placed in front of every statement:
\begin{lstlisting}
<label>: <statement>
\end{lstlisting}

As Java does not allow the use of \lstinline|goto|\index{goto}\footnote{… despite it is a reserved keyword: \autocite{ORACLE_DOC_LANGUAGE_SPECIFICATION:Keywords}.}, labels\index{label} are mainly used with the \lstinline|break|\index{break} and \lstinline|continue|\index{continue} statements. Consequently, labels will be assigned mostly to \lstinline|switch|\index{switch} (the \bdq{\texttt{case L:}} variant only), \lstinline|for|\index{for}, \lstinline|while|\index{while} and \lstinline|do-while|\index{while!do} statements.

\keeplisting{13}
But I found labels\index{label} particularly useful to give a name to an arbitary code block or a statement, in particular when the related code block is quite long and/or complex\footnote{Of course, you could also move this code into a private method. This is discussed in \tqfullvref{sec:InlineVsMethods}.}:
\begin{lstlisting}
TaxCalculation: {
    // Do the tax calculation here!
    …
}   //  TaxCalculation:

OptionalProcessing: if( doOptional() )
{
    // Do the optional handling here!
    …
}   //  OptionalProcessing:    
\end{lstlisting}

Of course, you can achieve something similar with comments, but the definition of an additional scope with the block statement is also useful, especially when dealing with lambdas\index{lambda} (see \tqfullvref{sec:Lambdas}.)

For a \lstinline|switch|\index{switch} statement, the \lstinline|break|\index{break} statement terminates a branch (refer to chapter \tqfullref{sec:TraditionalSwitch}).

In the context of loops, the \lstinline|break|\index{break} and the \lstinline|continue|\index{continue} statements are used to modify the execution flow of those loops\footnote{Some other coding standards recommend to completely avoid the use of \lstinline|break|\index{break} and \lstinline|continue|\index{continue} with loops. But in many cases, \lstinline|break|\index{break} and \lstinline|continue|\index{continue} can make the implementation of an algorithm easier and better understandable.}:

\begin{itemize}
	\item{\lstinline|break|\index{break} aborts the current loop cycle and terminates the loop; the execution flow continues with the first statement after the loop body for a \lstinline|for|\index{for} or an entry-controlled\index{entry-controlled loop} \lstinline|while|\index{while} loop, and after the \lstinline|while|\index{while!do} for an exit-controlled\index{exit-controlled loop} loop.}
	\item{\lstinline|continue|\index{continue} aborts the current loop cycle, too, but it continues with a check of the loop's termination condition. If this is not met yet, the loop body is executed again.}
\end{itemize}

I recommend to always use a meaningful label\index{label} together with \lstinline|break|\index{break} and \lstinline|continue|\index{continue}. This would make it more obvious what is intended with the \lstinline|break|\index{break} and \lstinline|continue|\index{continue}, and it assures that the right block will be terminated. If assigned, the label\index{label} should be used also as the end comment for a code block, as described in chapter \tqfullref{sec:TrailingOrEndOfLineComments} and chapter \tqfullref{sec:CommentsWhen}.

\keeplisting{11}
Here some samples on this topic; first a bad one:
\begin{lstlisting}
// BAD!!!
for( var line = 0; line < maxLines; ++line )
{
    for( var column = 0; column < maxColumns; ++column )
    {
        if( !isValid( field [line] [column] ) break; // What??
        …
    }
    processLine( field [line] );
}
\end{lstlisting}
Of course, the \jlsref defines unambigiously whether the \lstinline|break| in line~6 would terminate the inner or the outer loop, but do you now the respective chapters (\autocite{ORACLE_DOC_LANGUAGE_SPECIFICATION:Break} for \lstinline|break|\index{break} and \autocite{ORACLE_DOC_LANGUAGE_SPECIFICATION:Continue} for \lstinline|continue|\index{continue}) by heart?

\keeplisting{11}
The use of labels\index{label} clarifies the code's intent: 
\index{for}\index{break}
\begin{lstlisting}
// RECOMMENDED
LineLoop: for( var line = 0; line < maxLines; ++line )
{
    ColumnLoop: for( var column = 0; column < maxColumns; ++column )
    {
        if( !isValid( field [line] [column] ) break ColumnLoop;
        …
    } //  ColumnLoop:
    processLine( field [line] );
} // LineLoop:
\end{lstlisting}

\keeplisting{7}
\index{while}\index{continue}
\begin{lstlisting}
// RECOMMENDED
FilterLoop: while( input.hasNext() )
{
    final var value = input.next();
    if( !value.isValid() ) continue FilterLoop; // Skip invalid
    …
} // FilterLoop:
\end{lstlisting}

\keeplisting{12}
\index{switch}\index{switch!case}\index{switch!default}\index{break}\index{java.lang!IllegalArgumentException}
\begin{lstlisting}
// RECOMMENDED
DirectionSwitch: switch( direction )
{
    case LEFT: goLeft(); 
    	break DirectionSwitch;
    	
    case RIGHT: goRight(); 
    	break DirectionSwitch;
    	
    default:
        throw new IllegalArgumentException( direction.toString() );
} //  DirectionSwitch:
\end{lstlisting}

The name of the label\index{label} is repeated in a comment at the end of the code block. The colon at the end of the label name is mandatory for this kind of comment.

\subsubsection{Principles}
\begin{enumerate}[label=P\arabic*.]
\item{A label\index{label} can be used to give an arbitrary name to a code block.

It should be used with lengthy code blocks to simplify navigation inside the source code.

The label has to be repeated as a comment at the end of the code block; the colon is part of the label and may not be omitted.}
\item{There is no blank between the label itself and the colon, while the blank after the colon is required.

The statement follows the associated label on the same line.}
\item{It is strongly recommended to use labels together with the \lstinline|break|\index{break} and \lstinline|continue|\index{continue} statements.}
\end{enumerate}

%==============================================================================
% End of file!!
%==============================================================================
